\chapter{Plan de trabajo y cronograma}
\label{cap:plan-gantt}

El plan organiza el trabajo en cinco fases y conserva los hitos del calendario
académico de referencia para 2026/2027. El cronograma se divide en dos figuras
para mejorar la lectura: la figura~\ref{fig:gantt-fases-1-3} presenta el
cierre de requisitos, el diseño y las iteraciones de desarrollo; la
figura~\ref{fig:gantt-fases-4-5} presenta las pruebas, la memoria y el cierre.
Las fechas de las figuras son la planificación prevista. La ejecución conocida
a 31 de octubre de 2026 se resume en la tabla~\ref{tab:plan-previsto-real};
el registro del proyecto documenta estados y fechas de hitos, pero no conserva
un comienzo y cierre diario verificable para cada paquete de trabajo.

\begin{figure}[p]
\centering
\begin{ganttchart}[
  time slot format=isodate,
  x unit=0.070cm,
  y unit title=0.50cm,
  y unit chart=0.78cm,
  bar height=0.52,
  group height=0.45,
  milestone height=0.52,
  link tolerance=1pt,
  bar label font=\scriptsize,
  group label font=\scriptsize\bfseries,
  milestone label font=\scriptsize,
  bar label node/.append style={text width=5.6cm,align=right},
  group label node/.append style={text width=5.6cm,align=right},
  milestone label node/.append style={text width=5.6cm,align=right},
  hgrid,
  vgrid,
  title/.append style={fill=Color2!20},
  group/.append style={fill=Color2!45},
  bar/.append style={fill=Color1!75},
  milestone/.append style={fill=BackgroundColor,draw=ColorEtiqueta}
]{2026-09-01}{2026-12-31}
\gantttitlecalendar{year, month=shortname} \\
\ganttgroup{Fase 1. Cierre de requisitos y validación}{2026-09-01}{2026-10-15} \\
\ganttbar[name=rnf]{Mapeo RNF--ISO/IEC 25010 y justificación}{2026-09-01}{2026-09-25} \\
\ganttbar[name=uc]{Refactorización del diagrama de casos de uso}{2026-09-01}{2026-09-25} \\
\ganttbar[name=spec]{Corrección de especificaciones afectadas}{2026-09-15}{2026-10-14} \\
\ganttmilestone[name=mreq]{H1. Requisitos y diagrama validados}{2026-10-15} \\
\ganttnewline
\ganttnewline
\ganttgroup{Fase 2. Diseño}{2026-09-15}{2026-10-15} \\
\ganttbar[name=design]{Modelo de datos, multi-tenant, urna y registro}{2026-09-15}{2026-10-14} \\
\ganttmilestone[name=mdesign]{H2. Diseño técnico cerrado}{2026-10-15} \\
\ganttnewline
\ganttnewline
\ganttgroup{Fase 3. Desarrollo por iteraciones}{2026-10-16}{2026-12-31} \\
\ganttbar[name=i1]{Iteración 1. Organizaciones, roles y censo (RF-01--RF-03)}{2026-10-16}{2026-10-31} \\
\ganttbar[name=i2]{Iteración 2. Procesos y acceso (RF-06--RF-08)}{2026-11-01}{2026-11-15} \\
\ganttmilestone[name=mbase]{H3. Gestión y acceso funcionales}{2026-11-15} \\
\ganttbar[name=i3]{Iteración 3. Emisión y recibo (RF-09--RF-10)}{2026-11-16}{2026-11-30} \\
\ganttbar[name=i4]{Iteración 4. Auditoría, escrutinio y resultados}{2026-12-01}{2026-12-15} \\
\ganttbar[name=i5]{Iteración 5. Notificaciones y privacidad (RF-17--RF-18)}{2026-12-16}{2026-12-31} \\
\ganttlink{rnf}{mreq}
\ganttlink{uc}{mreq}
\ganttlink{spec}{mreq}
\ganttlink{design}{mdesign}
\ganttlink{mdesign}{i1}
\ganttlink{i1}{i2}
\ganttlink{i2}{mbase}
\ganttlink{mbase}{i3}
\ganttlink{i3}{i4}
\ganttlink{i4}{i5}
\end{ganttchart}

\vspace{0.5cm}
\begin{tabular}{lll}
  \tikz{\fill[Color2!45] (0,0) rectangle (0.45,0.22);} & Fase &
  \tikz{\fill[Color1!75] (0,0) rectangle (0.45,0.22);} Tarea \\
  \multicolumn{3}{c}{\tikz{\node[diamond,draw=ColorEtiqueta,fill=BackgroundColor,inner sep=2pt] {};}
  \quad Hito de entrega o validación}
\end{tabular}
\caption{Plan de trabajo, fases 1 a 3.}
\label{fig:gantt-fases-1-3}
\end{figure}

\begin{figure}[p]
\centering
\begin{ganttchart}[
  time slot format=isodate,
  x unit=0.070cm,
  y unit title=0.50cm,
  y unit chart=0.78cm,
  bar height=0.52,
  group height=0.45,
  milestone height=0.52,
  link tolerance=1pt,
  bar label font=\scriptsize,
  group label font=\scriptsize\bfseries,
  milestone label font=\scriptsize,
  bar label node/.append style={text width=5.6cm,align=right},
  group label node/.append style={text width=5.6cm,align=right},
  milestone label node/.append style={text width=5.6cm,align=right},
  hgrid,
  vgrid,
  title/.append style={fill=Color2!20},
  group/.append style={fill=Color2!45},
  bar/.append style={fill=Color1!75},
  milestone/.append style={fill=BackgroundColor,draw=ColorEtiqueta}
]{2026-10-16}{2027-02-15}
\gantttitlecalendar{year, month=shortname} \\
\ganttgroup{Fase 4. Pruebas}{2026-10-16}{2027-01-06} \\
\ganttbar[name=unit]{Pruebas unitarias incrementales por iteración}{2026-10-16}{2026-12-31} \\
\ganttbar[name=integ]{Integración de urna y registro de participación}{2026-12-01}{2026-12-20} \\
\ganttbar[name=perf]{Concurrencia (RNF-12) y carga (RNF-17)}{2026-12-16}{2027-01-03} \\
\ganttbar[name=ux]{Accesibilidad y usabilidad (RNF-1, RNF-3, RNF-4, RNF-6)}{2026-12-16}{2027-01-03} \\
\ganttmilestone[name=mtst]{H4. Evaluación y pruebas cerradas}{2027-01-06} \\
\ganttnewline
\ganttnewline
\ganttgroup{Fase 5. Validación y cierre}{2026-12-01}{2027-02-15} \\
\ganttbar[name=mem]{Memoria final, revisión del tutor y correcciones}{2026-12-01}{2027-01-17} \\
\ganttmilestone[name=mdir]{H5. Entrega al director}{2027-01-18} \\
\ganttbar[name=final]{Ajustes finales y entrega oficial}{2027-01-19}{2027-01-27} \\
\ganttmilestone[name=marati]{H6. Entrega oficial en ARATÍES}{2027-01-27} \\
\ganttmilestone[name=mvobo]{H7. Visto bueno del tutor}{2027-01-28} \\
\ganttbar[name=defprep]{Preparación de la defensa}{2027-01-15}{2027-02-07} \\
\ganttmilestone[name=mdef]{H8. Inicio del periodo de defensa}{2027-02-08} \\
\ganttlink{unit}{mtst}
\ganttlink{integ}{mtst}
\ganttlink{perf}{mtst}
\ganttlink{ux}{mtst}
\ganttlink{mem}{mdir}
\ganttlink{mtst}{mdir}
\ganttlink{mdir}{final}
\ganttlink{final}{marati}
\ganttlink{marati}{mvobo}
\ganttlink{mvobo}{mdef}
\ganttlink{defprep}{mdef}
\end{ganttchart}

\vspace{0.5cm}
\begin{tabular}{lll}
  \tikz{\fill[Color2!45] (0,0) rectangle (0.45,0.22);} & Fase &
  \tikz{\fill[Color1!75] (0,0) rectangle (0.45,0.22);} Tarea \\
  \multicolumn{3}{c}{\tikz{\node[diamond,draw=ColorEtiqueta,fill=BackgroundColor,inner sep=2pt] {};}
  \quad Hito de entrega o validación}
\end{tabular}
\caption{Plan de trabajo, fases 4 y 5.}
\label{fig:gantt-fases-4-5}
\end{figure}

\begin{longtable}{|p{0.24\textwidth}|p{0.25\textwidth}|p{0.16\textwidth}|p{0.23\textwidth}|}
\caption{Fechas, duraciones aproximadas y dependencias principales.}\label{tab:fechas-gantt}\\
\hline
\textbf{Tarea o hito} & \textbf{Inicio--fin} & \textbf{Duración} & \textbf{Dependencia} \\
\hline
\endfirsthead
\hline
\textbf{Tarea o hito} & \textbf{Inicio--fin} & \textbf{Duración} & \textbf{Dependencia} \\
\hline
\endhead
\hline
\endfoot
\hline
\endlastfoot
Cierre de requisitos y validación & 01/09/2026--15/10/2026 & 45 días & Inicio; revisión de requisitos \\
\hline
Diseño de datos y arquitectura & 15/09/2026--15/10/2026 & 31 días & En paralelo con el cierre de requisitos \\
\hline
Iteraciones 1--5 & 16/10/2026--31/12/2026 & 77 días & Diseño técnico y validación por módulos \\
\hline
Pruebas e integración & 16/10/2026--06/01/2027 & 83 días & Cada iteración; carga y accesibilidad tras disponer de módulos \\
\hline
Memoria y revisión del tutor & 01/12/2026--17/01/2027 & 48 días & Redacción progresiva; cierre tras H4 \\
\hline
Entrega al director & 18/01/2027 & Hito & Memoria y pruebas cerradas \\
\hline
Entrega oficial / visto bueno & 27/01/2027 / 28/01/2027 & Hitos & Entrega al director y correcciones finales \\
\hline
Preparación / defensa & 15/01/2027--08/02/2027 en adelante & 25 días hasta el inicio & Entrega oficial; defensa según convocatoria \\
\hline
\end{longtable}

Las pruebas unitarias avanzan durante las iteraciones. El diseño puede
comenzar mientras se cierra la validación de requisitos; los hitos H1--H8
marcan las revisiones y entregas. Las fechas internas son estimaciones y los
hitos oficiales deben confirmarse para la convocatoria correspondiente.

\section{Plan previsto y ejecución registrada}
\label{sec:plan-previsto-real}

La tabla compara las ventanas previstas para P0--P8 con el estado que figura
en el registro de progreso a 31 de octubre de 2026. «Hecho a la fecha de
corte» indica que el hito consta como completado, pero no implica que se
conozca su duración real. No se reconstruyen fechas de inicio a partir de
commits porque el historial no constituye un registro de esfuerzo.

\small
\setlength{\tabcolsep}{3pt}
\begin{longtable}{|p{0.06\textwidth}|p{0.22\textwidth}|p{0.24\textwidth}|p{0.34\textwidth}|}
\caption{Comparación entre planificación y estado de ejecución de P0--P8.}\label{tab:plan-previsto-real}\\
\hline
\textbf{Hito} & \textbf{Ventana prevista} & \textbf{Estado registrado a 31/10/2026} & \textbf{Desviación y observación} \\
\hline
\endfirsthead
\hline
\textbf{Hito} & \textbf{Ventana prevista} & \textbf{Estado registrado a 31/10/2026} & \textbf{Desviación y observación} \\
\hline
\endhead
\hline
\endfoot
\hline
\endlastfoot
P0 & 29/10--02/11/2026 & Hecho a la fecha de corte & La línea base está registrada; el cierre diario no está consignado. \\
\hline
P1 & 02/11--08/11/2026 & Hecho a la fecha de corte & La compilación inicial y la infraestructura constan como completadas antes de su ventana prevista. \\
\hline
P2 & 05/11--10/11/2026 & Hecho a la fecha de corte & Modelo, contratos y decisiones constan como completados antes de su ventana prevista. \\
\hline
P3 & 09/11--17/11/2026 & Hecho a la fecha de corte & Organizaciones, censo y base de interfaz constan como completados antes de su ventana prevista. \\
\hline
P4 & 16/11--30/11/2026 & Hecho a la fecha de corte & El flujo vertical RF-06--RF-10 consta como integrado antes de la iteración planificada. \\
\hline
P5 & 30/11--14/12/2026 & Hecho a la fecha de corte & Auditoría, escrutinio y resultados constan como completados antes de sus ventanas planificadas. \\
\hline
P6 & 14/12--31/12/2026 & Hecho a la fecha de corte & Consola y requisitos transversales constan como completados antes de su ventana planificada. \\
\hline
P7 & 01/01--06/01/2027 & En curso; 178 pruebas (119 de backend y 59 de frontend), 16 de ellas de integración sobre PostgreSQL & Se inició antes de la ventana; H1b y H1c ya se ejecutaron. E2E funcional, integración del resto de repositorios, revisión manual de accesibilidad y repetición de RNF-4 siguen pendientes. \\
\hline
P8 & 01/12/2026--27/01/2027 & En curso; diagnóstico P8.0 y borrador de capítulos & La revisión comenzó antes de la ventana prevista. Se recompiló un borrador y se hizo una inspección visual inicial; la revisión editorial final y la evaluación empírica de P7 siguen abiertas. \\
\hline
\end{longtable}
\setlength{\tabcolsep}{6pt}
\normalsize

La ejecución adelantó trabajo respecto a la secuencia prevista, pero no permite
concluir que todas las tareas se acortaran en la misma proporción. P7 y P8 se
solapan: la memoria incorpora la implementación existente y marca como
pendiente la evidencia de evaluación aún no recogida.
